\scanpage{606}\mspage{15}
% archive sheet number: 605

On a posé ici, bien sûr, pour tt filtre $\Phi$ sur $E$
\[
 k_{\Phi}=\varinjlim_{J\in\Phi}k^J .
\]

On \deleted{\unsure{s}}sait que \inserted{pour tte $k$-algèbre comm. $A$,}
$\mathrm{Hom}_M(\deleted{\unsure{k_{\Phi}}}A_{\unsure{\Phi|\deleted{\Phi}}},k)$ est en
corr. 1-1 avec la partie de $\mathrm{Max}(\deleted{\unsure{k}}A_{\Phi})$
formée des idéaux maximaux $m$ de $\deleted{\unsure{k}}A_{\Phi}$
tels que \deleted{$\unsure{A}$} $k\unsure{-|\to}A/m$ soit un iso.
Si $A\inserted{\text{ est de la forme }}\unsure{=}k_{\Phi}$, $\Phi$ un filtre
sur $E$, cet ensemble est en corr. bi-univoque \deleted{\unsure{a}}avec l'ens des
\inserted{$\pi'_0$-}ultrafiltres $\Psi$ sur $E$ plus fins que $\Phi$,
où $\pi'_0=\mathrm{Card}(k)$ (proposition 3, et cor. 3
\inserted{de prop. 2} page 7). \deleted{\unsure{Comme}} vu
$\pi'_0\geq\pi_0$, un tel ultrafiltre est trivial
\inserted{par hyp.}.

Donc on trouve
\[
\begin{aligned}
\mathrm{Hom}_M(k_{\Phi},k)&\simeq\text{ensemble des }\xi\in E\text{ tels que}\\
&\{\xi\}\in\deleted{\Phi}\bigcap_{A\in\Phi}A
\end{aligned}\tag*{(*)}
\]
Mais \deleted{\unsure{si}} \inserted{dire que} $\Phi$ est plus fin que
$\Phi_0$ signifie aussi que pour tt $\xi\in E$, on a
$\deleted{\unsure{A}}E\smallsetminus\{\xi\}\in\Phi$, i.e.
$\exists A\in\Phi$, \unsure{avec|où} $A\subset E-\{\xi\}$ i.e.
\deleted{ce qui implique $\{\xi\}\notin\Phi$ (}
\deleted{\unsure{donc}}\inserted{i.e.} $\{\xi\}\notin\bigcap_{A\in\Phi}A$.
Donc on a alors $\mathrm{Hom}_M(k_{\Phi},k)=\emptyset$.

\deleted{Notons maintenant que} La conclusion du cor. 1 s'exprime
par la relation
\[
 k_{\Phi}\notin\mathrm{Ob}\,\underline S\qquad\text{i.e.}\qquad
 k_{\Phi}\in\mathrm{Ob}\,\underline S',
\]
d'autre part on a par définition
\[
 k_{\Phi}=\varinjlim_{J\in\Phi}k_{\unsure{J|D}},
\]
